\chapter{Metodologia e architettura}
\label{cha:methodology}

\section{Architettura generale}
\label{sec:architecture}
L'obiettivo metodologico dell'esperimento consiste nella progettazione e valutazione di un sistema applicato ad attività di assemblaggio manuale nell'ambito di Predictive Process Monitoring (PPM), implementato per predire in anticipo il successo o il fallimento del montaggio associando o meno ai tradizionali registri di processo i segnali posturali dell'operatore. Si configura quindi come un task di \textit{outcome-prediction} (classificazione binaria) seguito da un modulo di spiegabilità dei modelli utilizzati.
\begin{figure}[htbp]
    \centering
    \inputtikz{images/source/pipeline_architecture.tex}
    \caption{Architettura della pipeline sperimentale}
    \label{fig:pipeline_architettura}
\end{figure}

Come illustrato nella Figura~\ref{fig:pipeline_architettura}, la pipeline si articola in 5 fasi principali:
\begin{enumerate}
    \item \textbf{Dataset}: caricamento dei log di assemblaggio ed estrazione delle feature nelle due configurazioni \texttt{no\_human} e \texttt{human}, con partizionamento in 10 fold stratificati per la Cross-Validation.
    \item \textbf{Bucketing ed encoding}: suddivisione e codifica delle tracce secondo tre diverse strategie di bucketing, con rappresentazione tabulare (2D) per i modelli classici e tensoriale (3D) per la rete neurale ricorrente (LSTM).
\item \textbf{Addestramento e ottimizzazione}: addestramento dei quattro classificatori (Decision Tree, Random Forest, XGBoost e LSTM) mediante Stratified Nested Cross-Validation e ottimizzazione bayesiana degli iperparametri (Hyperopt TPE) orientata alla massimizzazione del coefficiente MCC sul validation set.
    \item \textbf{Valutazione}: calcolo delle metriche di classificazione aggregate sui test set di tutti i 10 fold di Cross-Validation e identificazione del fold con MCC mediano come istanza rappresentativa. 
    \item \textbf{Spiegabilità}: attribuzione di feature importance tramite tecnica SHAP con confronto tra metodi di spiegabilità \textit{model-specific} (\texttt{TreeExplainer} per Decision Tree, Random Forest, XGBoost; \texttt{DeepExplainer} per LSTM) e agnostici (\texttt{KernelExplainer}).
\end{enumerate}

\section{Dataset}
\label{sec:dataset}
L'indagine sperimentale si innesta su una pipeline di preparazione dati precedentemente implementata, comprensiva delle fasi di \textit{event log extraction}, \textit{event log enrichment}, \textit{event log pre-processing} e \textit{prefix generation} (per maggiori dettagli si veda il lavoro svolto da Vian~\cite{ppm-human-centered-vian2024}).

Sono considerate due diverse configurazioni del dataset:

\paragraph{\texttt{no\_human} (baseline di processo).}
Include i dati dell'event log tradizionale. Rappresenta la configurazione di riferimento standard nell'ambito del Process Mining.

\paragraph{\texttt{human} (configurazione arricchita).}
Aggiunge alla baseline di processo dei segnali riferiti all'operatore umano. In particolare vengono calcolati i seguenti attributi:
\begin{itemize}
    \item \texttt{Action Execution Variants}: indica variazioni nell'esecuzione delle azioni di montaggio. È un attributo categorico definito in modo deterministico.
    \item \texttt{Action Quality}: specifica azioni eseguite non correttamente. È un attributo categorico definito in modo deterministico.
    \item \texttt{Focus of Attention}: indica se l'operatore è attento o meno al montaggio del pezzo. È un attributo binario, con valore \texttt{True} o \texttt{False}.
    \item \texttt{Pose Ergonomic Score}: indica una valutazione ergonomica di specifiche parti del corpo attraverso punteggi sulla scala RULA (utilizzata per descrivere il rischio di disordini muscolo-scheletrici). Sono incluse le metriche per braccio (\textit{upper arm}), avambraccio (\textit{lower arm}), polso (\textit{wrist}) e collo (\textit{neck}).
    \item \texttt{Movement Dynamics}: quantifica il livello di movimento dell'operatore. Comprende sia la media che la deviazione standard della posizione del tronco durante la stessa azione.
\end{itemize}

Le tracce reali del dataset sono già pre-suddivise secondo uno schema di Stratified Nested 10-Fold Cross-Validation, in cui ciascun fold preserva la proporzione delle classi e ripartisce i dati in \textbf{Training} (\textasciitilde 70\% delle tracce, per l'addestramento dei modelli), \textbf{Validation} (\textasciitilde 20\%, per l'ottimizzazione degli iperparametri) e \textbf{Testing} (10\%, riservato al calcolo delle metriche di performance finali).

\section{Bucketing ed encoding}
\label{sec:bucketing_encoding}
\subsection{Bucketing}
\label{subsec:bucketing}
L'obiettivo dell'indagine di \textit{outcome-prediction} è anticipare il risultato di una traccia di assemblaggio a partire dai prefissi incompleti, ottenuti riducendo l'intera esecuzione attraverso l'applicazione di differenti strategie di bucketing. Vian~\cite{ppm-human-centered-vian2024} indaga le seguenti tre strategie, riprodotte nella pipeline:
\paragraph{\texttt{PREFIX\_LENGTH}.}
Le tracce complete sono troncate a un numero fisso di eventi a partire dall'evento iniziale, con numero di eventi che varia da 5 a 25 con step unitari. Poiché tutti i prefissi dello stesso bucket hanno la stessa lunghezza, non vengono adottate strategie di padding.

\paragraph{\texttt{COMPONENT}.}
Le tracce sono troncate al completamento dell'assemblaggio di 1 dei 4 pannelli principali, risultando in 4 bucket di analisi. Per gestire diverse lunghezze nei prefissi dello stesso bucket viene adottata una strategia di NaN-padding.

\paragraph{\texttt{COMPONENT\_TIME\_WARPED}.}
Come \texttt{COMPONENT}, ma è adottata una strategia di padding basata su \textit{time-warping}. I prefissi sono espansi temporalmente (su azioni non atomiche) per adattarsi in termini di numero di eventi al più lungo prefisso osservato.

\begin{figure}[htbp]
    \centering
    %! TeX root = ../../main.tex
\documentclass[tikz,border=8pt,11pt]{standalone}
\usepackage[T1]{fontenc}
\usepackage{newtxtext,newtxmath}
\usepackage{xcolor}
\usepackage{tikz}
\usetikzlibrary{shapes.geometric, arrows.meta, positioning, fit, backgrounds, calc}

\begin{document}

% ==============================================================================
% DEFINIZIONE COLORI (PALETTE ACCADEMICA UNIFICATA)
% ==============================================================================
\definecolor{actA}{RGB}{252, 232, 131}          % A: Giallo caldo pastello
\definecolor{actB}{RGB}{242, 168, 178}          % B: Rosa corallo tenue (Azione non-atomica)
\definecolor{actC}{RGB}{188, 222, 178}          % C: Verde salvia
\definecolor{actD}{RGB}{164, 198, 238}          % D: Azzurro polvere (Azione non-atomica)
\definecolor{actNaN}{RGB}{235, 235, 235}        % NaN: Grigio neutro
\definecolor{actEmpty}{RGB}{255, 255, 255}      % Slot vuoto

\definecolor{hdr_bg}{RGB}{242, 245, 249}        % Sfondo intestazioni griglia
\definecolor{hdr_txt}{RGB}{25, 35, 50}          % Testo intestazioni griglia

\begin{tikzpicture}[
    font=\footnotesize,
    x=0.96cm,
    y=0.52cm,
    % Stile cella della griglia
    cell/.style={
        draw=black!80,
        line width=0.55pt,
        inner sep=0pt,
        minimum width=0.96cm,
        minimum height=0.52cm,
        font=\footnotesize,
        text=black,
        align=center
    },
    % Intestazione riga sinistra
    rowhdr/.style={
        draw=black!80,
        line width=0.55pt,
        fill=hdr_bg,
        inner sep=0pt,
        minimum width=2.9cm,
        minimum height=0.52cm,
        font=\footnotesize\bfseries,
        text=hdr_txt,
        align=left
    },
    % Intestazione colonna numerica (1..10)
    colhdr/.style={
        draw=black!80,
        line width=0.55pt,
        fill=hdr_bg,
        inner sep=0pt,
        minimum width=0.96cm,
        minimum height=0.52cm,
        font=\footnotesize\bfseries,
        text=hdr_txt,
        align=center
    },
    % Titolo sopra ogni tabella
    sectiontitle/.style={
        font=\bfseries\footnotesize,
        text=black,
        anchor=south west,
        align=left
    }
]

    % Coordinata X per allineamento a sinistra (bordo sinistro di rowhdr a x = -3.45)
    \def\leftX{-3.45}
    % Coordinata Y per il distacco del titolo dalla tabella
    \def\titleY{0.50}

    % ==========================================================================
    % 1. TRACCE COMPLETE DELL'EVENT LOG
    % ==========================================================================
    \node[sectiontitle] at (\leftX, \titleY) {Tracce complete dell'Event Log};

    % Header colonne
    \node[rowhdr] at (-2.0, 0) {\hspace{6pt}Numero evento};
    \foreach \i in {1,...,10} {
        \node[colhdr] at (\i - 1, 0) {\i};
    }

    % Trace 1
    \node[rowhdr] at (-2.0, -1) {\hspace{6pt}Traccia 1};
    \node[cell, fill=actA] at (0, -1) {A};
    \node[cell, fill=actB] at (1, -1) {B};
    \node[cell, fill=actA] at (2, -1) {A};
    \node[cell, fill=actC] at (3, -1) {C};
    \node[cell, fill=actD] at (4, -1) {D};
    \node[cell, fill=actB] at (5, -1) {B};
    \node[cell, fill=actA] at (6, -1) {A};
    \node[cell, fill=actD] at (7, -1) {D};
    \node[cell, fill=actEmpty] at (8, -1) {};
    \node[cell, fill=actEmpty] at (9, -1) {};

    % Trace 2
    \node[rowhdr] at (-2.0, -2) {\hspace{6pt}Traccia 2};
    \node[cell, fill=actA] at (0, -2) {A};
    \node[cell, fill=actC] at (1, -2) {C};
    \node[cell, fill=actA] at (2, -2) {A};
    \node[cell, fill=actA] at (3, -2) {A};
    \node[cell, fill=actB] at (4, -2) {B};
    \node[cell, fill=actD] at (5, -2) {D};
    \node[cell, fill=actD] at (6, -2) {D};
    \node[cell, fill=actC] at (7, -2) {C};
    \node[cell, fill=actA] at (8, -2) {A};
    \node[cell, fill=actA] at (9, -2) {A};

    % ==========================================================================
    % 2. PREFIX_LENGTH, lunghezza = 5
    % ==========================================================================
    \begin{scope}[yshift=-2.85cm]
        \node[sectiontitle] at (\leftX, \titleY) {\texttt{PREFIX\_LENGTH}, lunghezza $= 5$};

        % Header colonne
        \node[rowhdr] at (-2.0, 0) {\hspace{6pt}Numero evento};
        \foreach \i in {1,...,10} {
            \node[colhdr] at (\i - 1, 0) {\i};
        }

        % Prefix 1
        \node[rowhdr] at (-2.0, -1) {\hspace{6pt}Prefisso 1};
        \node[cell, fill=actA] at (0, -1) {A};
        \node[cell, fill=actB] at (1, -1) {B};
        \node[cell, fill=actA] at (2, -1) {A};
        \node[cell, fill=actC] at (3, -1) {C};
        \node[cell, fill=actD] at (4, -1) {D};
        \foreach \i in {5,...,9} {
            \node[cell, fill=actEmpty] at (\i, -1) {};
        }

        % Prefix 2
        \node[rowhdr] at (-2.0, -2) {\hspace{6pt}Prefisso 2};
        \node[cell, fill=actA] at (0, -2) {A};
        \node[cell, fill=actC] at (1, -2) {C};
        \node[cell, fill=actA] at (2, -2) {A};
        \node[cell, fill=actA] at (3, -2) {A};
        \node[cell, fill=actB] at (4, -2) {B};
        \foreach \i in {5,...,9} {
            \node[cell, fill=actEmpty] at (\i, -2) {};
        }
    \end{scope}

    % ==========================================================================
    % 3. COMPONENT, target = [A,B,C], componente = 2
    % ==========================================================================
    \begin{scope}[yshift=-5.70cm]
        \node[sectiontitle] at (\leftX, \titleY) {\texttt{COMPONENT}, target $= [\text{A},\text{B},\text{C}]$, componente $= 2$};

        % Header colonne
        \node[rowhdr] at (-2.0, 0) {\hspace{6pt}Numero evento};
        \foreach \i in {1,...,10} {
            \node[colhdr] at (\i - 1, 0) {\i};
        }

        % Prefix 1
        \node[rowhdr] at (-2.0, -1) {\hspace{6pt}Prefisso 1};
        \node[cell, fill=actA] at (0, -1) {A};
        \node[cell, fill=actB] at (1, -1) {B};
        \node[cell, fill=actA] at (2, -1) {A};
        \node[cell, fill=actC] at (3, -1) {C};
        \node[cell, fill=actD] at (4, -1) {D};
        \node[cell, fill=actB] at (5, -1) {B};
        \node[cell, fill=actNaN] at (6, -1) {\footnotesize\textit{NaN}};
        \node[cell, fill=actNaN] at (7, -1) {\footnotesize\textit{NaN}};
        \node[cell, fill=actEmpty] at (8, -1) {};
        \node[cell, fill=actEmpty] at (9, -1) {};

        % Prefix 2
        \node[rowhdr] at (-2.0, -2) {\hspace{6pt}Prefisso 2};
        \node[cell, fill=actA] at (0, -2) {A};
        \node[cell, fill=actC] at (1, -2) {C};
        \node[cell, fill=actA] at (2, -2) {A};
        \node[cell, fill=actA] at (3, -2) {A};
        \node[cell, fill=actB] at (4, -2) {B};
        \node[cell, fill=actD] at (5, -2) {D};
        \node[cell, fill=actD] at (6, -2) {D};
        \node[cell, fill=actC] at (7, -2) {C};
        \node[cell, fill=actEmpty] at (8, -2) {};
        \node[cell, fill=actEmpty] at (9, -2) {};
    \end{scope}

    % ==========================================================================
    % 4. COMPONENT_TIME_WARPED
    % ==========================================================================
    \begin{scope}[yshift=-8.55cm]
        \node[sectiontitle] at (\leftX, \titleY) {\texttt{COMPONENT\_TIME\_WARPED}, target $= [\text{A},\text{B},\text{C}]$, componente $= 2$, non-atomiche $= [\text{B},\text{D}]$};

        % Header colonne
        \node[rowhdr] at (-2.0, 0) {\hspace{6pt}Numero evento};
        \foreach \i in {1,...,10} {
            \node[colhdr] at (\i - 1, 0) {\i};
        }

        % Prefix 1 (D duplicato a pos 5 e 6, B duplicato a pos 7 e 8)
        \node[rowhdr] at (-2.0, -1) {\hspace{6pt}Prefisso 1};
        \node[cell, fill=actA] at (0, -1) {A};
        \node[cell, fill=actB] at (1, -1) {B};
        \node[cell, fill=actA] at (2, -1) {A};
        \node[cell, fill=actC] at (3, -1) {C};
        \node[cell, fill=actD] at (4, -1) {D};
        \node[cell, fill=actD] at (5, -1) {D};
        \node[cell, fill=actB] at (6, -1) {B};
        \node[cell, fill=actB] at (7, -1) {B};
        \node[cell, fill=actEmpty] at (8, -1) {};
        \node[cell, fill=actEmpty] at (9, -1) {};

        % Prefix 2
        \node[rowhdr] at (-2.0, -2) {\hspace{6pt}Prefisso 2};
        \node[cell, fill=actA] at (0, -2) {A};
        \node[cell, fill=actC] at (1, -2) {C};
        \node[cell, fill=actA] at (2, -2) {A};
        \node[cell, fill=actA] at (3, -2) {A};
        \node[cell, fill=actB] at (4, -2) {B};
        \node[cell, fill=actD] at (5, -2) {D};
        \node[cell, fill=actD] at (6, -2) {D};
        \node[cell, fill=actC] at (7, -2) {C};
        \node[cell, fill=actEmpty] at (8, -2) {};
        \node[cell, fill=actEmpty] at (9, -2) {};
    \end{scope}

\end{tikzpicture}

\end{document}

    \caption{Esempio di generazione dei prefissi utilizzando le tre strategie descritte sopra}
    \label{fig:encoding_strategies}
\end{figure}

\subsection{Encoding}
\label{subsec:encoding}
Successivamente vengono adottate le adeguate strategie di \textit{encoding} per mappare i prefissi nei formati utilizzabili dai classificatori considerati. Diversamente da Vian~\cite{ppm-human-centered-vian2024} che utilizza esclusivamente Decision Tree, nella pipeline descritta i prefissi sono stati trasformati in
\begin{enumerate*}[label=(\roman*)]
    \item formato tabulare per Decision Tree, Random Forest e XGBoost
    \item formato tensoriale per LSTM
\end{enumerate*}.

In entrambi i casi i dati si sviluppano considerando le seguenti 3 dimensioni:
\begin{itemize}
    \item $N$: numero di istanze analizzate appartenenti allo stesso bucket
    \item $T$: numero di passi temporali considerati
    \item $D$: cardinalità delle feature di base
\end{itemize}


\paragraph{Formato tabulare.}
La rappresentazione tabulare di dimensione $(N \times (T \cdot D))$ viene impiegata per i modelli basati su alberi (Decision Tree, Random Forest e XGBoost). Come osservato da Vian, poiché i prefissi contengono attributi sia categorici che continui, va adottata una strategia ibrida: i dati continui vengono mantenuti nel loro valore originario, i dati categorici sono tradotti secondo una logica di \textit{one-hot encoding}. La dimensione temporale è rappresentata orizzontalmente su più attributi secondo la convenzione \texttt{<feature>\_<step t>}.
\begin{figure}[H]
    \centering
    %! TeX root = ../../main.tex
\documentclass[tikz,border=10pt,11pt]{standalone}
\usepackage[T1]{fontenc}
\usepackage{newtxtext,newtxmath}
\usepackage{xcolor}
\usepackage{tikz}
\usetikzlibrary{shapes.geometric, arrows.meta, positioning, fit, backgrounds, calc, decorations.pathreplacing}

\begin{document}

% ==============================================================================
% DEFINIZIONE COLORI (COERENTE CON ENCODING_STRATEGIES E SETTINGS)
% ==============================================================================
\definecolor{hdr_bg}{RGB}{242, 245, 249}        % Sfondo intestazioni griglia
\definecolor{hdr_txt}{RGB}{25, 35, 50}          % Testo intestazioni griglia
\definecolor{pad_bg}{RGB}{248, 249, 250}        % Sfondo tenue per celle di padding
\definecolor{pad_txt}{RGB}{100, 100, 100}       % Testo grigio per padding

\begin{tikzpicture}[
    font=\footnotesize,
    % Stile cella intestazione (colhdr)
    colhdr/.style={
        draw=black!80,
        line width=0.55pt,
        fill=hdr_bg,
        inner sep=0pt,
        minimum height=0.58cm,
        font=\footnotesize\bfseries,
        text=hdr_txt,
        align=center
    },
    % Stile cella dati standard
    cell/.style={
        draw=black!80,
        line width=0.55pt,
        fill=white,
        inner sep=0pt,
        minimum height=0.58cm,
        font=\footnotesize,
        text=black,
        align=center
    },
    % Stile cella per padding
    padcell/.style={
        draw=black!80,
        line width=0.55pt,
        fill=pad_bg,
        inner sep=0pt,
        minimum height=0.58cm,
        font=\footnotesize\itshape,
        text=pad_txt,
        align=center
    }
]

    % ==========================================================================
    % PARAMETRI GEOMETRICI: LARGHEZZE E CENTRI COLONNE PERFETTAMENTE ALLINEATI
    % ==========================================================================
    % Colonne:
    % 1: Traccia      [ 0.00,  1.75] -> w = 1.75, cx = 0.88
    % 2: Label        [ 1.75,  2.65] -> w = 0.90, cx = 2.20
    % 3: act_1        [ 2.65,  3.65] -> w = 1.00, cx = 3.15
    % 4: rula_1       [ 3.65,  4.75] -> w = 1.10, cx = 4.20
    % 5: act_2        [ 4.75,  5.75] -> w = 1.00, cx = 5.25
    % 6: rula_2       [ 5.75,  6.85] -> w = 1.10, cx = 6.30
    % 7: act_3        [ 6.85,  8.10] -> w = 1.25, cx = 7.47
    % 8: rula_3       [ 8.10,  9.45] -> w = 1.35, cx = 8.78

    \def\wColA{1.75cm} \def\xColA{0.88}
    \def\wColB{0.90cm} \def\xColB{2.20}
    \def\wColC{1.00cm} \def\xColC{3.15}
    \def\wColD{1.10cm} \def\xColD{4.20}
    \def\wColE{1.00cm} \def\xColE{5.25}
    \def\wColF{1.10cm} \def\xColF{6.30}
    \def\wColG{1.25cm} \def\xColG{7.47}
    \def\wColH{1.35cm} \def\xColH{8.78}

    % Quote Y dei centri riga (altezza cella = 0.58cm)
    \def\yHdr{0.0}
    \def\yRowA{-0.58}
    \def\yRowB{-1.16}

    % ==========================================================================
    % RIGA 0: INTESTAZIONE (HEADER)
    % ==========================================================================
    \node[colhdr, minimum width=\wColA] at (\xColA, \yHdr) {Traccia};
    \node[colhdr, minimum width=\wColB] at (\xColB, \yHdr) {Label};
    \node[colhdr, minimum width=\wColC] at (\xColC, \yHdr) {\textbf{act\_1}};
    \node[colhdr, minimum width=\wColD] at (\xColD, \yHdr) {\textbf{rula\_1}};
    \node[colhdr, minimum width=\wColE] at (\xColE, \yHdr) {\textbf{act\_2}};
    \node[colhdr, minimum width=\wColF] at (\xColF, \yHdr) {\textbf{rula\_2}};
    \node[colhdr, minimum width=\wColG] at (\xColG, \yHdr) {\textbf{act\_3}};
    \node[colhdr, minimum width=\wColH] at (\xColH, \yHdr) {\textbf{rula\_3}};

    % ==========================================================================
    % RIGA 1: CASSETTO 01 (Label = 1, traccia regolare completa)
    % ==========================================================================
    \node[cell, minimum width=\wColA] at (\xColA, \yRowA) {Cassetto 01};
    \node[cell, minimum width=\wColB] at (\xColB, \yRowA) {1};
    \node[cell, minimum width=\wColC] at (\xColC, \yRowA) {2};
    \node[cell, minimum width=\wColD] at (\xColD, \yRowA) {3.5};
    \node[cell, minimum width=\wColE] at (\xColE, \yRowA) {4};
    \node[cell, minimum width=\wColF] at (\xColF, \yRowA) {5.0};
    \node[cell, minimum width=\wColG] at (\xColG, \yRowA) {1};
    \node[cell, minimum width=\wColH] at (\xColH, \yRowA) {2.0};

    % ==========================================================================
    % RIGA 2: CASSETTO 05 (Label = 0, anomala con padding su t=3)
    % ==========================================================================
    \node[cell,    minimum width=\wColA] at (\xColA, \yRowB) {Cassetto 05};
    \node[cell,    minimum width=\wColB] at (\xColB, \yRowB) {0};
    \node[cell,    minimum width=\wColC] at (\xColC, \yRowB) {2};
    \node[cell,    minimum width=\wColD] at (\xColD, \yRowB) {4.0};
    \node[cell,    minimum width=\wColE] at (\xColE, \yRowB) {3};
    \node[cell,    minimum width=\wColF] at (\xColF, \yRowB) {6.5};
    \node[padcell, minimum width=\wColG] at (\xColG, \yRowB) {0 \textit{(pad)}};
    \node[padcell, minimum width=\wColH] at (\xColH, \yRowB) {0.0 \textit{(pad)}};

    % ==========================================================================
    % GRAFFE DIMENSIONALI: N A DESTRA, T x D SOTTO
    % ==========================================================================
    % Graffa a destra: dimensione N (righe delle istanze/tracce)
    \draw[black!85, line width=0.55pt, decorate, decoration={brace, amplitude=4pt}]
        (9.57, -0.29) -- (9.57, -1.45)
        node[midway, right=5pt, font=\small] {$N$};

    % Graffa in basso: dimensione T x D (colonne di feature concatenate)
    \draw[black!85, line width=0.55pt, decorate, decoration={brace, mirror, amplitude=4pt}]
        (2.65, -1.57) -- (9.45, -1.57)
        node[midway, below=5pt, font=\small] {$T \cdot D$};

\end{tikzpicture}

\end{document}

    \caption{Esempio di prefisso trasformato in formato tabulare}
    \label{fig:tabular_format}
\end{figure}

\paragraph{Formato tensoriale.}
Per l'architettura LSTM i prefissi trasformati in formato tabulare vengono mappati in tensori tridimensionali di forma $(N, T, D)$. La dinamica temporale del processo viene così preservata lungo la sequenza, consentendo alla rete ricorrente di elaborare gli eventi in ordine cronologico.
\begin{figure}[H]
    \centering
    %! TeX root = ../../main.tex
\documentclass[tikz,border=10pt,11pt]{standalone}
\usepackage[T1]{fontenc}
\usepackage{newtxtext,newtxmath}
\usepackage{xcolor}
\usepackage{tikz}
\usetikzlibrary{shapes.geometric, arrows.meta, positioning, fit, backgrounds, calc, decorations.pathreplacing}

\begin{document}

% ==============================================================================
% DEFINIZIONE COLORI (ARMONIZZATA CON TABULAR FORMAT ED ENCODING STRATEGIES)
% ==============================================================================
\definecolor{hdr_bg}{RGB}{242, 245, 249}
\definecolor{hdr_txt}{RGB}{25, 35, 50}
\definecolor{pad_bg}{RGB}{248, 249, 250}
\definecolor{pad_txt}{RGB}{100, 100, 100}
\definecolor{depth_line}{RGB}{140, 155, 175}

\begin{tikzpicture}[
    font=\footnotesize,
    >=Stealth,
    colhdr/.style={
        draw=black!80,
        line width=0.55pt,
        fill=hdr_bg,
        inner sep=0pt,
        minimum height=0.58cm,
        font=\footnotesize\bfseries,
        text=hdr_txt,
        align=center
    },
    cell/.style={
        draw=black!80,
        line width=0.55pt,
        fill=white,
        inner sep=0pt,
        minimum height=0.58cm,
        font=\footnotesize,
        text=black,
        align=center
    },
    padcell/.style={
        draw=black!80,
        line width=0.55pt,
        fill=pad_bg,
        inner sep=0pt,
        minimum height=0.58cm,
        font=\footnotesize\itshape,
        text=pad_txt,
        align=center
    }
]

    % Dimensioni colonne e righe (uguali a tabular_format.tex)
    \def\wStep{0.80cm}
    \def\wAct{1.30cm}
    \def\wRula{1.45cm}
    \def\hRow{0.58cm}

    % Dimensioni complessive tabella: W = 0.80 + 1.30 + 1.45 = 3.55cm
    % Altezza complessiva: 4 * 0.58 = 2.32cm (da y = -1.16 a 1.16)
    \def\W{3.55}
    \def\Htop{1.16}
    \def\Hbot{-1.16}

    % Offset prospettico lungo asse N (profondita)
    \def\dx{3.85}
    \def\dy{1.15}

    % ==========================================================================
    % 1. LINEE TRATTEGGIATE PROSPETTICHE DI SFONDO
    % ==========================================================================
    % Vertice superiore sinistro
    \draw[depth_line, dashed, line width=0.55pt] (0, \Htop) -- ++(\dx, \dy);
    % Vertice inferiore sinistro
    \draw[depth_line, dashed, line width=0.55pt] (0, \Hbot) -- ++(\dx, \dy);
    % Vertice inferiore destro
    \draw[depth_line, dashed, line width=0.55pt] (\W, \Hbot) -- ++(\dx, \dy);

    % ==========================================================================
    % 2. FETTA POSTERIORE: Cassetto 05 (completamente visibile)
    % ==========================================================================
    \begin{scope}[shift={(\dx, \dy)}]
        % Titolo fetta posteriore
        \node[font=\footnotesize\bfseries, text=hdr_txt, anchor=south west] 
            at (0, \Htop + 0.10) {Cassetto 05};

        % Intestazioni colonne
        \node[colhdr, minimum width=\wStep] at (0.400, 0.87) {$t$};
        \node[colhdr, minimum width=\wAct]  at (1.450, 0.87) {\textbf{act}};
        \node[colhdr, minimum width=\wRula] at (2.825, 0.87) {\textbf{rula}};

        % Riga 1: t=1
        \node[colhdr, minimum width=\wStep] at (0.400, 0.29) {$1$};
        \node[cell, minimum width=\wAct]    at (1.450, 0.29) {2};
        \node[cell, minimum width=\wRula]   at (2.825, 0.29) {4.0};

        % Riga 2: t=2
        \node[colhdr, minimum width=\wStep] at (0.400, -0.29) {$2$};
        \node[cell, minimum width=\wAct]    at (1.450, -0.29) {3};
        \node[cell, minimum width=\wRula]   at (2.825, -0.29) {6.5};

        % Riga 3: t=3 (padding)
        \node[colhdr, minimum width=\wStep]  at (0.400, -0.87) {$3$};
        \node[padcell, minimum width=\wAct]  at (1.450, -0.87) {0 \textit{(pad)}};
        \node[padcell, minimum width=\wRula] at (2.825, -0.87) {0.0 \textit{(pad)}};
    \end{scope}

    % ==========================================================================
    % 3. FETTA ANTERIORE: Cassetto 01
    % ==========================================================================
    \begin{scope}[shift={(0, 0)}]
        % Titolo fetta anteriore con sfondo bianco per non farsi tagliare da linee
        \node[font=\footnotesize\bfseries, text=hdr_txt, anchor=south west, fill=white, inner sep=2pt] 
            at (0, \Htop + 0.10) {Cassetto 01};

        % Intestazioni colonne
        \node[colhdr, minimum width=\wStep] at (0.400, 0.87) {$t$};
        \node[colhdr, minimum width=\wAct]  at (1.450, 0.87) {\textbf{act}};
        \node[colhdr, minimum width=\wRula] at (2.825, 0.87) {\textbf{rula}};

        % Riga 1: t=1
        \node[colhdr, minimum width=\wStep] at (0.400, 0.29) {$1$};
        \node[cell, minimum width=\wAct]    at (1.450, 0.29) {2};
        \node[cell, minimum width=\wRula]   at (2.825, 0.29) {3.5};

        % Riga 2: t=2
        \node[colhdr, minimum width=\wStep] at (0.400, -0.29) {$2$};
        \node[cell, minimum width=\wAct]    at (1.450, -0.29) {4};
        \node[cell, minimum width=\wRula]   at (2.825, -0.29) {5.0};

        % Riga 3: t=3
        \node[colhdr, minimum width=\wStep] at (0.400, -0.87) {$3$};
        \node[cell, minimum width=\wAct]    at (1.450, -0.87) {1};
        \node[cell, minimum width=\wRula]   at (2.825, -0.87) {2.0};
    \end{scope}

    % ==========================================================================
    % 4. ELEMENTI IN PRIMO PIANO: LINEA SUPERIORE DESTRA E GRAFFE DIMENSIONALI
    % ==========================================================================
    % Linea tratteggiata del vertice in alto a destra (in primo piano)
    \draw[depth_line, dashed, line width=0.55pt] (\W, \Htop) -- ++(\dx, \dy);

    % Dimensione T (Passi temporali - graffa con mirror a sinistra della fetta frontale)
    \draw[black!85, line width=0.55pt, decorate, decoration={brace, mirror, amplitude=4pt}]
        (-0.10, 0.58) -- (-0.10, \Hbot)
        node[midway, left=5pt, font=\small] {$T$};

    % Dimensione D (Feature - graffa sotto la fetta frontale, su act e rula)
    \draw[black!85, line width=0.55pt, decorate, decoration={brace, mirror, amplitude=4pt}]
        (0.80, \Hbot - 0.12) -- (\W, \Hbot - 0.12)
        node[midway, below=5pt, font=\small] {$D$};

    % Dimensione N (Graffa sotto la linea che collega i vertici in basso a destra)
    \draw[black!85, line width=0.55pt, decorate, decoration={brace, mirror, amplitude=4pt, raise=4pt}]
        (\W, \Hbot) -- ++(\dx, \dy)
        node[midway, below right=6pt, font=\small] {$N$};

\end{tikzpicture}

\end{document}

    \caption{Esempio di prefisso trasformato in formato tensoriale}
    \label{fig:tensor_format}
\end{figure}

\section{Addestramento e ottimizzazione}
\label{sec:addestramento_ottimizzazione}
Successivamente la pipeline prosegue con l'addestramento e l'ottimizzazione di più classificatori. I passi principali sono:
\begin{enumerate*}[label=(\roman*)]
    \item la scelta dei modelli da implementare
    \item l'ottimizzazione dei modelli utilizzando training e validation set
\end{enumerate*}.


\subsection{Panoramica dei classificatori}
Sono stati scelti modelli comunemente utilizzati per i task di Predictive Process Monitoring, indagando l'utilizzo di classificatori ad albero e Long Short-Term Memory (LSTM). Diversamente da Vian~\cite{ppm-human-centered-vian2024}, il cui obiettivo è la sola valutazione dell'impatto delle feature umane, lo scopo dell'esperimento è analizzare la qualità delle predizioni con modelli più robusti ma più complessi nella loro interpretazione.

\begin{enumerate}
    \item \textbf{Decision Tree}: già impiegato da Vian~\cite{ppm-human-centered-vian2024}, funge da baseline di riferimento.
    \item \textbf{Random Forest}: introdotto come modello \textit{ensemble} di tipo \textit{bagging}, per valutare un miglioramento nelle predizioni con l'aggregazione di più alberi.
    \item \textbf{XGBoost}: modello \textit{ensemble} scelto per valutare l'efficacia del \textit{gradient boosting} sui dati tabulari rispetto all'applicazione del \textit{bagging} delle Random Forest.
    \item \textbf{LSTM}: transizione al \textit{deep learning}, scelto per indagare la dinamica temporale in modo sequenziale e non codificata in formato tabulare.
\end{enumerate}

\subsection{Procedura di addestramento e ottimizzazione degli iperparametri}
Come conseguenza della procedura di addestramento è addestrato e salvato un modello per ogni
\begin{enumerate*}[label=(\roman*)]
    \item sottogruppo di dataset,
    \item classificatore e
    \item strategia di generazione dei prefissi
\end{enumerate*}.

Come nell'architettura di Vian~\cite{ppm-human-centered-vian2024}, la pipeline segue uno schema di Stratified Nested 10-Fold Cross-Validation:
\begin{itemize}
    \item Inner loop: il modello viene addestrato sui dati di \textit{training}, mentre i dati di \textit{validation} sono utilizzati nell'algoritmo di ricerca bayesiana per identificare la configurazione di iperparametri che massimizza il coefficiente MCC. Nel caso di XGBoost e LSTM sono calcolate anche delle metriche di \textit{early-stopping} (rispettivamente il numero ottimale di alberi e il numero di epoche ideali).
    \item Outer loop: selezionati i migliori iperparametri viene addestrato un modello finale sulla concatenazione dei dati di \textit{training} e \textit{validation}, poi valutato sul \textit{testing set}. Tale approccio permette di estendere la configurazione trovata su un dataset più numeroso e migliorare la generalizzazione del modello. Nel caso di XGBoost e LSTM il modello candidato viene eseguito per il numero di alberi o epoche identificato precedentemente.
\end{itemize}

\subsection{Spazio di ricerca degli iperparametri}
Per ogni classificatore sono stati considerati gli spazi di ricerca riportati nella Tabella~\ref{tab:search_spaces}:

\begin{table}[H]
    \centering
    \small
    \renewcommand{\arraystretch}{1.15}
    \begin{tabular}{ll}
        \toprule
        \textbf{Iperparametro} & \textbf{Distribuzione / Spazio di ricerca} \\
        \midrule
        \multicolumn{2}{l}{\textbf{Decision Tree} (\textit{budget}: 500 iterazioni)} \\
        \midrule
        \texttt{max\_depth} & $\{2, 4, 6, 8, 10, 12, 14, 16, 18, 20, 24, 30, 40, 50, \text{None}\}$ \\
        \texttt{min\_samples\_split} & $\{2, 4, 6, 8, 10, 13, 16, 20, 30, 40\}$ \\
        \texttt{min\_samples\_leaf} & $\{2, 4, 6, 8, 10, 13, 16, 20, 30, 40\}$ \\
        \texttt{criterion} & $\{\text{gini}, \text{entropy}\}$ \\
        \texttt{ccp\_alpha} & $\mathcal{U}(0,\, 0.05)$ \\
        \midrule
        \multicolumn{2}{l}{\textbf{Random Forest} (\textit{budget}: 500 iterazioni)} \\
        \midrule
        \texttt{n\_estimators} & $\{100, 200, 300, 400, 500\}$ \\
        \texttt{max\_depth} & $\{2, 4, 6, 8\}$ \\
        \texttt{min\_samples\_split} & $\{2, 4, 6, 8, 10, 13, 16, 20, 30, 40\}$ \\
        \texttt{min\_samples\_leaf} & $\{2, 4, 6, 8, 10, 13, 16, 20, 30, 40\}$ \\
        \texttt{criterion} & $\{\text{gini}, \text{entropy}\}$ \\
        \midrule
        \multicolumn{2}{l}{\textbf{XGBoost} (\textit{budget}: 500 iterazioni)} \\
        \midrule
        \texttt{n\_estimators} & $\{100, 200, 300, 500\}$ \\
        \texttt{max\_depth} & $\{2, 3\}$ \\
        \texttt{learning\_rate} & $\log\mathcal{U}(-5,\, -1) \approx [0.007,\, 0.368]$ \\
        \texttt{subsample} & $\mathcal{U}(0.5,\, 1.0)$ \\
        \texttt{colsample\_bytree} & $\mathcal{U}(0.3,\, 0.8)$ \\
        \texttt{gamma} & $\mathcal{U}(0,\, 5)$ \\
        \texttt{reg\_alpha} & $\log\mathcal{U}(-10,\, 2) \approx [4.5\times 10^{-5},\, 7.39]$ \\
        \texttt{reg\_lambda} & $\log\mathcal{U}(-10,\, 3) \approx [4.5\times 10^{-5},\, 20.09]$ \\
        \midrule
        \multicolumn{2}{l}{\textbf{LSTM} (\textit{budget}: 100 iterazioni)} \\
        \midrule
        \texttt{hidden\_dim} & $\{64, 128, 256\}$ \\
        \texttt{num\_layers} & $\{1, 2, 3\}$ \\
        \texttt{batch\_size} & $\{16, 32, 64\}$ \\
        \texttt{dropout} & $\mathcal{U}(0.1,\, 0.5)$ \\
        \texttt{lr} & $\log\mathcal{U}(-9.21,\, -4.6) \approx [10^{-4},\, 10^{-2}]$ \\
        \texttt{weight\_decay} & $\log\mathcal{U}(-13.8,\, -4.6) \approx [10^{-6},\, 10^{-2}]$ \\
        \bottomrule
    \end{tabular}
    \caption{Spazi di ricerca degli iperparametri per ciascun classificatore}
    \label{tab:search_spaces}
\end{table}



\section{Valutazione}
\label{sec:valutazione}
Lo schema di Stratified Nested 10-Fold Cross-Validation prevede per ogni iterazione 1 fold di \textit{test} e i rimanenti 9 fold suddivisi tra \textit{training} e \textit{validation}. In particolare, le metriche di valutazione finali sono calcolate mediando i risultati sui \textit{test set} di tutti e 10 i fold, mentre il modello candidato (utilizzato nei passi successivi) corrisponde al classificatore il cui coefficiente MCC è il valore mediano tra i 10 calcolati, ordinati in ordine crescente.

In continuità con lo studio di Vian~\cite{ppm-human-centered-vian2024}, la valutazione delle predizioni si concentra su tre metriche derivate dalla matrice di confusione binaria: Precision, Recall e il Matthews Correlation Coefficient (MCC).

Precision e Recall variano da 0 a 1 e indicano rispettivamente l'affidabilità delle predizioni positive e la sensibilità del modello nell'identificare le istanze positive reali.
\begin{equation}
    \begin{gathered}
        \text{Precision} = \frac{TP}{TP + FP}, \qquad
        \text{Recall} = \frac{TP}{TP + FN} \\[1.5ex]
    \end{gathered}
    \label{eq:precision_recall_metrics}
\end{equation}

Il coefficiente MCC assume valori compresi tra -1 (predizioni completamente errate), 0 (predizioni casuali) e +1 (predizioni corrette), fornendo una stima robusta che combina tutti i valori della matrice di confusione.
\begin{equation}
    \text{MCC} = \frac{TP \times TN - FP \times FN}{\sqrt{(TP + FP)(TP + FN)(TN + FP)(TN + FN)}}
    \label{eq:mcc_metric}
\end{equation}

\section{Spiegabilità}
\label{sec:spiegabilita}

\section{Stack tecnologico}
\label{sec:stack}